\begin{abstract}
    We present the \pokeagent Challenge, a large-scale benchmark for decision-making research built on Pokémon's multi-agent battle system and expansive role-playing game (RPG) environment. Partial observability, game-theoretic reasoning, and long-horizon planning remain open problems for frontier AI, yet few benchmarks stress all three simultaneously under realistic conditions. \pokeagent targets these limitations at scale through two complementary tracks: our \textit{Battling Track}, which calls for strategic reasoning and generalization under partial observability in competitive Pokémon battles, and our \textit{Speedrunning Track}, which requires long-horizon planning and sequential decision-making in the Pokémon RPG. Our Battling Track supplies a dataset of 20M+ battle trajectories alongside a suite of heuristic, RL, and LLM-based baselines capable of high-level competitive play. Our Speedrunning Track provides the first standardized evaluation framework for RPG speedrunning, including an open-source multi-agent orchestration system that enables modular, reproducible comparisons of harness-based LLM approaches. Our NeurIPS 2025 competition validates both the quality of our resources and the research community's interest in Pokémon, with more than 100 teams competing across both tracks and winning solutions detailed in our paper. Participant submissions and our own baselines reveal considerable gaps between generalist (LLM), specialist (RL), and elite human performance. Analysis against the BenchPress evaluation matrix shows that Pokémon battling is nearly orthogonal to standard LLM benchmarks, measuring capabilities not captured by existing evaluation suites and positioning Pokémon as an unsolved benchmark that can drive RL and LLM research forward. We transition from the NeurIPS 2025 competition to a living benchmark by releasing a live leaderboard for Battling and self-contained evaluation for Speedrunning at \url{https://pokeagentchallenge.com}.
    \end{abstract}
    
    
